% secciones/03_desarrollo.tex
\section{Desarrollo e Integración del Proyecto}\label{sec:desarrollo}

Esta sección documenta la implementación técnica de la solución integral para \textit{ML Lab}, detallando los componentes desarrollados en el backend de Spring Boot, el microservicio analítico de FastAPI y la interfaz de usuario reactiva en Angular.

\subsection{Capa de Backend (Java Spring Boot)}

La capa del servidor se diseñó para centralizar el flujo del negocio y actuar como pasarela intermedia. La persistencia de datos históricos se administra con Spring Data JPA. En el \autoref{lst:experiment_entity} se muestra la entidad mapeada a PostgreSQL para registrar las ejecuciones.

\lstinputlisting[
    language=Java,
    caption={Fragmento de la Entidad JPA Experiment para la persistencia en Supabase},
    label={lst:experiment_entity}
]{experiment_entity_fragmento.java}

El backend procesa las solicitudes de analítica y delega la ejecución al microservicio Python mediante peticiones HTTP síncronas. Adicionalmente, cuenta con endpoints dedicados para descargar fichas técnicas del entrenamiento en PDF. En el \autoref{lst:pdf_controller} se detalla la exposición de este endpoint.

\lstinputlisting[
    language=Java,
    caption={Fragmento del Controlador REST para la descarga de PDFs},
    label={lst:pdf_controller}
]{pdf_controller_fragmento.java}

La creación física del PDF se realiza a través de Apache PDFBox. El servicio define el documento de tamaño estándar A4 y dibuja dinámicamente formas decorativas y cadenas de texto sobre un canvas de coordenadas como se visualiza en el \autoref{lst:pdf_service}.

\lstinputlisting[
    language=Java,
    caption={Fragmento del Servicio para la construcción de reportes con Apache PDFBox},
    label={lst:pdf_service}
]{pdf_service_fragmento.java}

\subsection{Microservicio Analítico (Python FastAPI)}

El servidor en Python FastAPI expone de forma directa la lógica de modelado matemático. El archivo ejecutor principal configura las rutas agrupadas para cada algoritmo, tal como se muestra en el \autoref{lst:app_fastapi}.

\lstinputlisting[
    language=Python,
    caption={Estructura principal del servidor en FastAPI},
    label={lst:app_fastapi}
]{app_fastapi_fragmento.py}

La lógica operativa de entrenamiento y predicción está implementada utilizando la biblioteca de scikit-learn. El \autoref{lst:knn_model} ilustra el proceso para el clasificador de vecinos más cercanos (KNN).

\lstinputlisting[
    language=Python,
    caption={Modelo clasificador KNN con cálculo de exactitud y predicción en scikit-learn},
    label={lst:knn_model}
]{knn_model_fragmento.py}

% [CAPTURA PENDIENTE: Swagger UI mostrando todos los endpoints]
% \begin{figure}[H]
% \centering
% \includegraphics[width=0.8\textwidth,alt={Interfaz interactiva de Swagger UI listando endpoints de regresion lineal y clasificacion KNN}]{swagger_endpoints.png}
% \caption{Interfaz interactiva de documentación Swagger UI}
% \label{fig:swagger_endpoints}
% \end{figure}

\subsection{Capa de Frontend (Angular 18+)}

La interfaz web consume los recursos mediante servicios reactivos inyectables. El \autoref{lst:ml_service} ilustra el servicio encargado de despachar los flujos analíticos enviando los payloads correspondientes.

\lstinputlisting[
    language=JavaScript,
    caption={Servicio Angular reactivo para el consumo de la API},
    label={lst:ml_service}
]{ml_service_fragmento.ts}

Por último, el componente encapsula la lógica de interacción del formulario capturando la entrada en cadenas de texto formateadas en JSON y enviando las variables procesadas a través del servicio, como se detalla en el \autoref{lst:ml_knn_component}.

\lstinputlisting[
    language=JavaScript,
    caption={Componente controlador de la interfaz de usuario en Angular},
    label={lst:ml_knn_component}
]{ml_knn_component_fragmento.ts}
